\documentclass[10pt,a4paper]{amsart}
\usepackage[margin=19mm]{geometry}
\usepackage{amsmath,amssymb,mathtools}
\usepackage[scr=rsfs,cal=euler]{mathalfa}
\usepackage{xfrac}
\usepackage[unicode,hidelinks,bookmarks=false]{hyperref}
\newcommand{\ie}{i.e.\ }
\newcommand{\cf}{cf.\ }
\newcommand{\resp}{respectively\ }
\newcommand{\Subsection}{\subsection}
\providecommand{\nobreakdash}{\nobreak\hbox{-}}
\setlength{\emergencystretch}{2em}
\multlinegap0pt
\usepackage{amssymb}
\usepackage{xcolor}
\usepackage{tikz}
\usepackage{stackengine}
\usepackage{enumerate}
\usepackage[shortlabels]{enumitem}
\newtheorem{theorem}{Theorem}
\newcommand{\M}{\mathbf{M}}
\newcommand{\R}{\mathbf{R}}
\newcommand{\V}{\mathbf{V}}
\newcommand{\W}{\mathbf{W}}
\newcommand{\boxr}{\Box\mathsf{R}}
\newcommand{\calc}{\mathsf{CSGL}}
\newcommand{\ct}{\mathsf{ct}}
\newcommand{\cte}{\mathsf{E}}
\newcommand{\ctl}{\mathsf{L}}
\newcommand{\ctv}{\mathsf{V}}
\newcommand{\false}{\mathtt{false}}
\newcommand{\fval}{\false} %\mathtt{F}}
\newcommand{\id}{\mathsf{id}}
\newcommand{\iffi}{\textit{iff} }
\newcommand{\imp}{\rightarrow}
\newcommand{\impl}{{\rightarrow}\mathsf{L}}
\newcommand{\lab}{\mathsf{Lab}}
\newcommand{\ndrule}{\Box_\mathsf{DB}}
\renewcommand{\phi}{\varphi}
\newcommand{\prove}{\mathtt{prove}}
\newcommand{\ru}{\mathsf{r}}
\newcommand{\sar}{\vdash}
\newcommand{\scomp}{\odot}
\newcommand{\set}[1]{\{#1\}}
\newcommand{\tel}{\mathcal{T}}
\newcommand{\tseq}{T} %\mathscr{T}}
\title{Optimizing Proof-Search via Linearization for Gödel-Löb Logic with Tree-Hypersequents}
\author{Tim S. Lyon, Omar Taher}
\date{Source: arXiv:2606.03484v2 (CC BY 4.0)}
\begin{document}
\maketitle

\noindent\emph{Source and licence.} Optimizing Proof-Search via Linearization for Gödel-Löb Logic with Tree-Hypersequents. Version arXiv:2606.03484v2 (CC BY 4.0), distributed by arXiv under the Creative Commons Attribution 4.0 International licence, http://creativecommons.org/licenses/by/4.0/, as recorded in the arXiv licence metadata for this exact version and shown on the abstract page https://arxiv.org/abs/2606.03484v2. Full licence notice: \texttt{licenses/arxiv-2606.03484-CC-BY-4.0.txt}.\par\smallskip

\noindent\textit{Editorial scope.} Counted unit: the article's theorem lem:false-gives-refute (source0.tex lines 1277--1279). The article's complete original proof of it is delivered with it (source0.tex lines 1281--1320, one \texttt{proof} environment of 40 lines). No mathematics is written, repaired or re-derived: the statement and the proof are the article's own lines, in the article's own notation, and no retained line is substituted or re-flowed.  No further source-local context is needed: every cross-reference of every retained line resolves inside the two delivered spans.  Nothing was omitted from any retained span, so the package carries no omission record.  No retained line contains a citation, so the wrapper carries no bibliography and no bibliography entry.  No retained line contains a figure environment, an image-inclusion command or drawing code, so no image is delivered and the package declares no asset dependency.  Editorial formatting only, in addition: an AMS \texttt{amsart} preamble that restates the article's own declarations and macros used by the retained lines, namely the environments \texttt{theorem} on separate counters, together with the standard packages those lines need.  The retained environments print in this document as Theorem 1 in order of appearance, whereas the article prints the same items with the numbering of its own full document.  The complete native source of the article as distributed in the arXiv e-print for this exact version is bundled unchanged as \texttt{sources/201978/source0.tex}.
\begin{theorem}\label{lem:false-gives-refute}
If $\prove( \sar x : \phi) = \false$, then a model $\M = (\W,\R,\V)$ can be extracted from the corresponding computation tree such that $\M \not\models \phi$.
\end{theorem}

\begin{proof} Suppose $\prove(\sar x : \phi) = \false$ and let $\ct(x : \phi) = (\ctv,\cte,\ctl)$ be the corresponding computation tree. We prune the computation tree to obtain a structure $(\ctv',\cte')$ from which we can extract a counter-model for $\phi$. Let us define $(\ctv',\cte')$ as follows:
\begin{description}

\item[$(1)$] Let $( \sar x : \phi) \in \ctv'$ and observe that $\ctl( \sar x : \phi) = \fval$ by assumption;

\item[$(2)$] If $\ctl(\tseq) = \fval$ and $\tseq \in \ctv'$ concludes a unary rule $ \ru \in \calc \setminus \set{\boxr}$ in $\ct(x : \phi)$ with $\tseq' \in \ctv$ the premise, then $\tseq' \in \ctv'$ and $(\tseq,\tseq') \in \cte'$;

\item[$(3)$] If $\ctl(\tseq) = \fval$ and $\tseq \in \ctv'$ is the conclusion of $\impl$ in $\ct(x : \phi)$ with $\tseq_{1}, \tseq_{2} \in \ctv$ the premises, then for some $i \in \set{1,2}$, $\ctl(\tseq_{i}) = \fval$, so for exactly one such $i$, we let $\tseq_{i} \in \ctv'$ and $(\tseq,\tseq_{i}) \in \cte'$;

\item[$(4)$] If $\ctl(\tseq) = \fval$ and $\tseq \in \ctv'$ concludes $\ndrule$ in $\ct(x : \phi)$ with $\tseq_{1}, \ldots, \tseq_{n} \in \ctv$ the premises, then we let $\tseq_{1}, \ldots, \tseq_{n} \in \ctv'$ and $(\tseq,\tseq_{1}), \ldots, (\tseq,\tseq_{n}) \in \cte'$.

\end{description}
Observe that the structure $(\ctv',\cte')$ is obtained by starting at the root and taking the downward closure of sequents labeled with $\false$, with the exception that only a single premise of an $\impl$ application is retained (i.e., one premise labeled with $\false$ is retained while the other is ignored, regardless of its label). Hence, any branching that occurs in $(\ctv',\cte')$ is due to a $\ndrule$ rule application.

Let $\tseq_{1}, \ldots, \tseq_{n} \in \ctv'$ be all stable leaves in the structure $(\ctv',\cte')$, and define $\tseq^{*} := \tseq_{1} \scomp \cdots \scomp \tseq_{n}$. We let $\tseq_{i} := \tel_{i}, \Gamma_{i} \sar \Delta_{i}$ for $i \in [n]$, $\tel := \bigcup_{i \in [n]} \tel_{i}$, $\Gamma := \bigcup_{i \in [n]} \Gamma_{i}$, and $\Delta := \bigcup_{i \in [n]} \Delta_{i}$, so that $\tel^{*} = \tel, \Gamma \sar \Delta$. Recall that $\tseq_{1}, \ldots, \tseq_{n}$ must be line sequents. We now define the model $\M = (\W,\R,\V)$ such that (1) $\W := \lab(\tseq^{*})$, (2) $(y,z) \in \R$ \iffi there exist $u_{1}, \ldots, u_{k} \in \lab(\tseq^{*})$ such that $yRu_{1}, \ldots, u_{k}Rz \in \tel$, and (3) $y \in \V(p)$ \iffi $y : p \in \Gamma$. We now prove that $\M$ is indeed a model. 

First, since $\sar x : \phi$ was the input to proof-search, we know that $x \in \W$, and so, $\W \neq \emptyset$. Second, by construction, we know that $\tel$ is a finite tree, meaning, $\R$ is a finite transitively-closed tree. Hence, $\R$ is both transitive and conversely well-founded. Last, observe that $\V$ is well-defined. 

To finish the proof, we prove the following two claims by a mutual induction on the length of $\phi$ and $\psi$, for all $y \in \lab(\tseq^{*})$: (i) if $y : \phi \in \Gamma$, then $\M, y \models \phi$ and (ii) if $y : \psi \in \Delta$, then $\M, y \not\models \psi$.
\begin{description}

\item[$\phi = p.$] If $y : p \in \Gamma$, then by the definition of $\V$, we know that $y \in \V(p)$, and so, $\M, y \models p$.

\item[$\psi = p.$] Let $y : p \in \Delta$. Assume for a contradiction that $y : p \in \Gamma$ as well. Then, it must be the case that for some $i \neq j \in [n]$, $y : p \in \Gamma_{i}$ and $y : p \in \Delta_{j}$. Observe that if $i = j$, then $\tseq_{i}$ would not be saturated, contradicting our assumption that $\tseq_{i}$ is stable. $\tseq_{i}$ and $\tseq_{j}$ must occur along different branches of $(\ctv',\cte')$ as they are distinct leaves. Let $\tseq$ be the closest common ancestor of $\tseq_{i}$ and $\tseq_{j}$. Since $(\ctv',\cte')$ is a tree, such an ancestor must exist, and by what was said above, it must be the conclusion of a $\ndrule$ application. All labels shared by $\tseq_{i}$ and $\tseq_{j}$ must occur in $\tseq$ by construction because after $\ndrule$ is applied, all labels introduced will be fresh and pairwise distinct; consequently, $y \in \lab(\tseq)$. By the definition of $\prove$, we know that all rules in $(\ctv',\cte')$ will be end-active, meaning, after $\ndrule$ is applied bottom-up to $\tseq$, $y : p$ cannot be introduced along the branch to $\tseq_{i}$ or $\tseq_{j}$. Therefore, $y : p$ must occur in the antecedent and consequent of $\tseq$, which contradicts our assumption that $\ndrule$ was applied bottom-up to $\tseq$; $\ndrule$ is only applied to saturated sequents and in this case $\tseq$ would not satisfy condition ($\id$). It follows that $y : p \not\in \Gamma$, meaning $y \not\in \V(p)$, and so, $\M, y \not\models p$.

\item[$\phi = \bot.$] Observe that $y : \bot \not\in \Gamma$ since then some $\tseq_{i}$ would not be saturated, contrary to our assumption. Hence, claim (i) vacuously holds.

\item[$\psi = \bot.$] If $y : \bot \in \Delta$, then claim (ii) vacuously holds because $\M, y \not\models \bot$ by definition.

\item[$\phi = \chi \imp \theta.$] If $y : \chi \imp \theta \in \Gamma$, then there exists some $\tseq_{i}$ with $y : \chi \imp \theta \in \Gamma_{i}$. Since $\tseq_{i}$ is saturated, we know that either $y : \chi \in \Delta_{i} \subseteq \Delta$ or $y : \theta \in \Gamma_{i} \subseteq \Gamma$. By IH, either $\M,y \not\models \chi$ or $\M, y \models \theta$. Either way, $\M, y \models \chi \imp \theta$.

\item[$\psi = \chi \imp \theta.$] If $y : \chi \imp \theta \in \Delta$, then there exists some $\tseq_{i}$ with $y : \chi \imp \theta \in \Delta_{i}$. Since $\tseq_{i}$ is saturated, we know that $y : \chi \in \Gamma_{i} \subseteq \Gamma$ and $y : \theta \in \Delta_{i} \subseteq \Delta$. By IH, $\M,y \models \chi$ and $\M, y \not\models \theta$. Hence, $\M, y \not\models \chi \imp \theta$.

\item[$\phi = \Box \chi.$] Suppose $y : \Box \chi \in \Gamma$. Let $z \in \W$ with $(y,z) \in \R$. Then, there exist $u_{1}, \ldots, u_{k} \in \lab(\tseq^{*})$ such that $yRu_{1}, \ldots, u_{k}Rz \in \tel$ by definition. Since $\tel$ is the composition of $n$ line sequents, we know that some $i \in [n]$ exists such that $yRu_{1}, \ldots, u_{k}Rz \in \tel_{i}$ for $\tseq_{i} = \tel_{i}, \Gamma_{i} \sar \Delta_{i}$. As $\tseq_{i}$ is saturated, we know that $z : \chi \in \Gamma_{i} \subseteq \Gamma$. By IH, $\M, z \models \chi$, meaning, $\M, y \models \Box \chi$ since $z$ was arbitrary.

\item[$\psi = \Box \chi.$] Suppose $y : \Box \chi \in \Delta$. Then, there exists some stable line sequent $\tseq_{i} = \tel_{i}, \Gamma_{i} \sar \Delta_{i}$ such that $y : \Box \chi \in \Delta_{i}$. Since $\tseq_{i}$ is stable and $y : \Box \chi \in \Delta_{i}$, it cannot be the case that $y$ is a leaf; otherwise, $\tseq_{i}$ would not be stable. Hence, there must exist a $z \in \lab(\tseq^{*})$ such that $yRz \in \tel_{i}$. Let us consider the $\ndrule$ application on the path from $\tseq_{i}$ to the root of $(\ctv',\cte')$ which introduced $yRz$ with $z$ fresh, and let $\tseq_{i}' = \tel_{i}', \Gamma_{i}' \sar \Delta_{i}'$ be the conclusion of $\ndrule$. Since $y : \Box \chi \in \Delta_{i}$, by inspection of the rules applied during proof-search, one will find that $y : \Box \chi \in \Delta_{i}'$. Furthermore, observe that $\ctl(\tseq_{i}') = \fval$, meaning, every premise of the $\ndrule$ application will be labeled with $\fval$ as well. Thus, there will exist some premise $\tseq_{j}' = \tel_{j}', \Gamma_{j}' \sar \Delta_{j}'$ such that $yRu \in \tel_{j}'$ and $u : \chi \in \Delta_{j}'$. By the definition of $(\ctv',\cte')$, there will exist a stable tree sequent $\tseq_{j} = \tel_{j}, \Gamma_{j} \sar \Delta_{j}$ that is a leaf in $(\ctv',\cte')$ above $\tseq_{j}'$ such that $yRu \in \tel_{j}$ and $u : \chi \in \Delta_{j}$. Consequently, $yRu \in \tel$ and $u : \chi \in \Delta$, so by the definition of $\M$ and IH, we know that there exists a $u \in \W$ such that $(y,u) \in \R$ and $\M, u \not\models \chi$. This implies that $\M, y \not\models \Box \chi$.

\end{description}
Observe that $x : \phi \in \Delta$; by claim (ii), $\M \not\models \phi$.
\end{proof}

\end{document}
